\paragraph{The number of exceptional coordinates.}
\label{par:row-restricted}
At component \(c\) of the target, use the lattice
\(L=N_{c-1}/XN_{c-1}\) of Section~\ref{par:row-metric}.
Identify \(L\otimes_O{\rm Frac}(O)\)
with the crystalline constant space and pull every labeled filtration
back to it. At a plus embedding, \(C_\tau={\rm Fil}^{b_\tau}_\tau\)
has codimension at most \(m\). At a minus embedding,
\(B_\tau={\rm Fil}^{1-b_\tau}_\tau\) has dimension at most \(m\).
These inequalities follow by counting the central jumps, including
their multiplicities, on the two sides of the respective cuts.
The filtered spaces for distinct roots belonging to the same embedding
give the same test on the crystalline constants.

\begin{proposition}[A bounded set of exceptional coordinates]\label{prop:row-exceptional-coordinates}
Assume \(\operatorname{codim} C_\tau\le m\) at the plus
embeddings and \(\dim B_\tau\le m\) at the minus embeddings,
with at most \(e\) embeddings in total. There is a direct summand \(B\) with complement \(L_{\rm main}\) such that
\[
 L=L_{\rm main}\oplus B,\quad L_{\rm main}\subseteq\bigcap_\tau C_\tau,\quad B_\tau\subseteq B\otimes K\quad(\tau\text{ minus}),\quad \operatorname{rk}_O B\le em,
\]
where \(K\) is the fraction field and the intersection is over the plus embeddings.
The image \(I_B\) of the invariant residual coordinates defined
below is free of rank at most \((e+1)\operatorname{rk}_O B\).
\end{proposition}

\begin{proof}
\proofstep{Splitting off the exceptional directions.}
There is an integral splitting \(L=L_{\rm main}\oplus L_{\rm err}\)
such that \(L_{\rm main}\) lies in every \(C_\tau\) and
\(L_{\rm err}\otimes{\rm Frac}(O)\) contains every \(B_\tau\), with
\({\rm rank}\,L_{\rm err}\le em\). Let \(K=\operatorname{Frac}(O)\), \(C=\bigcap_{\tau,+}C_\tau\),
and write \(L=(C\cap L)\oplus D\) using saturation. Define
\[
  L_{\rm err}=L\cap\Big(D_K+\sum_{\tau,-}B_\tau\Big),
  \qquad
  \operatorname{rk}L_{\rm err}\le \operatorname{codim} C
                          +\sum_{\tau,-}\dim B_\tau\le em .
\]
Choose \(L_{\rm main}\) complementary to \(L_{\rm err}\cap C\)
in \(C\cap L\). Then \(L=L_{\rm main}\oplus L_{\rm err}\)
since \(D\subseteq L_{\rm err}\). All splittings used here exist over the DVR because
a saturated submodule of a finite free module has free quotient.
In particular the projection to \(L_{\rm main}\) has integral entries
and kills all the \(B_\tau\).

\proofstep{The polynomial that separates the main coordinates.}
Consider a model coordinate \(z\in E^{-H}A\otimes L\) with
\[
 (\gamma-1)z\in XN^{\rm mod},
\]
for example \(\beta(w)\) with \((\gamma-1)w\in XN_c\). At a primitive root the coefficient of
\(t^i\) in \(z\), for \(i\ne0\), belongs to \({\rm Fil}^{-i}_\tau\).
Indeed the coefficient of its difference is multiplied by \(u^i-1\),
nonzero in the fraction field, and multiplication by \(X\), a unit
at this root, preserves the filtration inequalities. On projecting
to \(L_{\rm main}\) it follows that the coefficients vanish in degrees
\(i<-b_\tau-S\) at plus embeddings, and in nonzero degrees
\(i<b_\tau\) at minus embeddings. In the latter case the coefficient
is in \(B_\tau\), which the projection kills.

Let \(\Omega\) be the union over these embeddings of the integer
intervals \([-b_\tau-S,-b_\tau)\) at plus and
\([b_\tau,b_\tau+S)\) at minus. The exceptional degrees determine the polynomial
\[
 P(Z)=\prod_{i\in\Omega}(Z-u^i).
\]
The operator \(P(\gamma)\) preserves \(E^{-H}A\) since \(\gamma E/E\)
is a unit. In the Laurent parameter at that root, write the projection
to the main summand as \({\rm pr}_{\rm main}z=\sum_i z_i t^i\).
Since this projection commutes with \(\gamma\),
\[
 [t^i]{\rm pr}_{\rm main}P(\gamma)z=P(u^i)z_i,\qquad
 P(u^i)=0\quad(i\in\Omega).
\]
Below the indicated intervals the projected coefficients already vanish,
with the possible exception of the constant at a minus embedding.
The polynomial removes the degrees inside the intervals. Write the main part of \(P(\gamma)z\) in the integral
basis obtained above as \(\sum f_j\otimes l_j\).

At plus points \(f_j\) has order at least \(-b_\tau\); at minus
points
\[
 f_j=x_{\tau,j}+O(t^{b_\tau+S})
\]
for a constant
\(x_{\tau,j}\) in the coefficient field (with the value at the
chosen root understood). Also \((\gamma-1)f_j\in X E^{-H}A\).
The latter inclusion follows by commuting projection and the polynomial
with \(\gamma\) and using the stability of \(XE^{-H}A\).
Each simple tensor is invariant modulo \(XN^{\rm mod}\).
For its difference divided by \(X\), the required filtration at
plus points follows from \(l_j\in{\rm Fil}^{b_\tau}\) and the lower
order just proved. At minus points the constant cancels and all
of \(L\) is in \({\rm Fil}^{-b_\tau-S}\).
Thus the residual part is also invariant modulo \(XN^{\rm mod}\).

Finally
\[
 P(\gamma)z\equiv P(1)z\bmod XN^{\rm mod},
\]
by stability
and \((Z-1)\mid(P(Z)-P(1))\).
The scalar cost to replace \(P(1)\) by a power of \(p\) is at most
\[
             \kappa=e\{S+v_p(S!)+\lfloor\log_p H\rfloor\}.
\]
Each interval has \(S\) nonzero integers of magnitude at most \(H\).
For one interval \(I\), the principal-unit formula therefore gives
\[
 \sum_{i\in I}v(1-u^i)
 = |I|+\sum_{a=1}^{\lfloor\log_p H\rfloor}
                  \#\{i\in I:p^a\mid i\}
 \le S+v_p(S!)+\lfloor\log_p H\rfloor .
\]
There are at most \(e\) intervals; overlaps can only reduce
the sum over their union. An empty union has cost zero.

\proofstep{Counting the residual directions.}
We need the number of integral residual classes, including the
coefficient from the ring variable. For \(B=L_{\rm err}\), consider
the image \(I_B\) in \((E^{-H}A\otimes L)/XN^{\rm mod}\) of those
\(z\in E^{-H}A\otimes B\) with \((\gamma-1)z\in XN^{\rm mod}\).
The two assertions to prove are torsion-freeness and the bound
\[
 \operatorname{rank}_O I_B\le(e+1)\operatorname{rank}_O B.
\]
Over the coefficient DVR they imply the asserted freeness once finiteness is known.
The ambient quotient is finite over \(O\): \(N^{\rm mod}\) contains
\(E^HA\otimes L\) by the jump bounds, and \(XE^{2H}\) is a unit times
a distinguished polynomial.

The ambient quotient is also torsion-free. If \(a z\in
XN^{\rm mod}\) for \(0\ne a\in O\), the constant coefficient of the
integral numerator of \(a z\) at \(X=0\) vanishes. Thus \(z/X\) still
has an integral numerator. At every primitive root the coefficient test for \(az/X\) lies in a
filtration space over \(K=\operatorname{Frac}(O)\). For each degree \(i\),
\[
  a[t^i](z/X)\in {\rm Fil}^{-i}_\tau
      \quad\Longrightarrow\quad [t^i](z/X)\in {\rm Fil}^{-i}_\tau .
\]
Together with integrality at zero this gives \(z/X\in N^{\rm mod}\).
The ambient quotient is consequently finite torsion-free over the DVR.
Its submodule \(I_B\) is free.

For the rank, use the finite filtration starting with \(E^{-H}A\otimes B\).
Multiply successively by \(X-r_j\), \(2H\) times for each of the
\(e\) roots \(r_j\) of \(E\), and finally by \(X\). The last term
is contained in \(XN^{\rm mod}\). Every grade has rank at most
\({\rm rank}\,B\), by division with remainder at its root, including
the last root zero. The filtration is \(\gamma\)-stable. At \(r_j\)
the local parameter is a unit times the cyclotomic logarithm. To see
this, write
\[
 X=\pi^s v(\pi)
\]
with \(v\) a unit.

At
\(\pi=\zeta-1\) its derivative is
\(\pi^{s-1}(sv(\pi)+\pi v'(\pi))\ne0\), since \(s\) is a \(p\)-adic
unit. The \(p-1\) simple roots in the ordinary parameter give \(e\)
simple roots in the invariant ring: the tame orbits have size \(s\)
and \(E(X(\pi))\) is a unit times the cyclotomic polynomial at
\(1+\pi\). The identical multiplicity calculation also shows that
different orbits give distinct roots. Alternatively \(E\) has
degree \(e\) and its derivative is nonzero at each image, and division
over each fibre of \(X-r_j\) in the ordinary disc has degree \(s\).

On the grade with preceding multiplicity \(a\) at \(r_j\), \(\gamma\)
acts by \(u^{a-H}\); factors coming from other roots are units there.
Since \(u\) has infinite order, the invariant-grade test at this root is
\[
 u^{a-H}=1\quad\Longleftrightarrow\quad a=H.
\]
The
grade at zero contributes at most one more copy of \(B\).
Over the fraction field filter the image in the quotient by
\(XN^{\rm mod}\). Its graded pieces are quotients of the original
grades, and a fixed vector has a fixed leading term. The action in the graded factors at \(r_j\) gives
\[
 \operatorname{rk}_O I_B
 \le \sum_{j=1}^{e}\sum_{\substack{0\le a<2H\\u^{a-H}=1}}
           \operatorname{rk}_O B+\operatorname{rk}_O B
 \le (e+1)\operatorname{rk}_O B\le e(e+1)m .
\]
Thus an \(O\)-basis of \(I_B\) furnishes a single choice of
representatives at this component for all target classes. Taking a different basis for each input
would give no bound on the number of exceptional generators.
\label{cor:row-coordinates-1}
\end{proof}
